% =========================================================
% Legacy-Archiv: legacy-oscillators.tex
% =========================================================
% Diese Datei enthaelt zusammengesetzte Schwinger-Pics aus mechlib v1. Sie wird
% vom zentralen Loader NICHT automatisch geladen und ist kein Bestandteil der
% neuen atomaren mech/*-API.
%
% Zweck der Datei:
%   - alte Abbildungen nachvollziehbar halten,
%   - Geometrien und Parameterlisten bei der Migration als Referenz bewahren,
%   - schrittweise Neuimplementierung mit mass/cart, spring und damper erlauben.
%
% WICHTIG: Die alten Pics verwenden historische Stil- und Pic-Namen wie
% "wagen", "feder", "daempfer" und "x_cos". Diese Abhaengigkeiten muessen im
% alten Dokument vorhanden sein. Neue Abbildungen sollen nicht auf diese API
% aufbauen. Die vorhandenen Abschnittskommentare dokumentieren jeweils Pic-Name
% und Parameterreihenfolge der einzelnen Legacy-Geometrien.
% =========================================================
\tikzset{
	pics/einmassenschwinger_anregung/.style args={#1,#2,#3,#4,#5}{
		code={
			\pgfmathsetmacro{\mh}{0.8*#1}
			\pgfmathsetmacro{\mb}{0.5*#2}
			\pgfmathsetmacro{\mr}{0.1*#1}
			\pgfmathsetmacro{\mf}{2*\mr+0.25*\mh}
			\pgfmathsetmacro{\md}{2*\mr+0.75*\mh}
				% Schwinger
			\pic[rotate=180] at (1.2*\mb,2*\mr+1.3*\mh){x_cos={0.4cm,0.2cm,ucbblue}};
			\node[ucbblue,anchor=south] at (1.2*\mb,2*\mr+1.3*\mh) {$x$};
			\pic at (\mb,0){wagen={\mh,\mb,\mr,$m$}};
			\pic {feder={0,\mb,\mf,$c$}};
			\pic {daempfer={0,\mb,\md,$d$}};
			\draw (1.5*\mb,2*\mr+1.*\mh)--+(0,2*\mr);
			\draw[red,-latex] (1.5*\mb,3*\mr+1.*\mh)--+(0.8,0)node[pos=.5,anchor=south]{$F$};
			% Wand
			\draw[](0,0) --++ (0,\mh+2*\mr);
			\fill[pattern={Lines[angle=45, distance={1pt}]},pattern color=black](0,0) rectangle (-3pt,\mh+2*\mr);
		}
	}
}
% ------------------ Einmassenschwinger ------------------
% Pic: einmassenschwinger
% Parameter: #1 Höhe, #2 Breite,#3 Name der Masse, #4 Federsteifigkeit, #5 Dämpfung
\tikzset{
	pics/einmassenschwinger_anregung_ungedaempft/.style args={#1,#2,#3,#4,#5}{
		code={
			\pgfmathsetmacro{\mh}{0.8*#1}
			\pgfmathsetmacro{\mb}{0.5*#2}
			\pgfmathsetmacro{\mr}{0.1*#1}
			\pgfmathsetmacro{\mf}{2*\mr+0.5*\mh}
			% Schwinger
			\pic[rotate=180] at (1.2*\mb,2*\mr+1.3*\mh){x_cos={0.4cm,0.2cm,ucbblue}};
			\node[ucbblue,anchor=south] at (1.2*\mb,2*\mr+1.3*\mh) {$x$};
			\pic at (\mb,0){wagen={\mh,\mb,\mr,$m$}};
			\pic {feder={0,\mb,\mf,$c$}};
			\draw (1.5*\mb,2*\mr+1.*\mh)--+(0,2*\mr);
			\draw[red,-latex] (1.5*\mb,3*\mr+1.*\mh)--+(0.8,0)node[pos=.5,anchor=south]{$F$};
			% Wand
			\draw[](0,0) --++ (0,\mh+2*\mr);
			\fill[pattern={Lines[angle=45, distance={1pt}]},pattern color=black](0,0) rectangle (-3pt,\mh+2*\mr);
		}
	}
}
% ------------------ Einmassenschwinger ------------------
% Pic: einmassenschwinger
% Parameter: #1 Höhe, #2 Breite,#3 Name der Masse, #4 Federsteifigkeit, #5 Dämpfung
\tikzset{
	pics/einmassenschwinger/.style args={#1,#2,#3,#4,#5}{
		code={
			\pgfmathsetmacro{\mh}{0.8*#1}
			\pgfmathsetmacro{\mb}{0.5*#2}
			\pgfmathsetmacro{\mr}{0.1*#1}
			\pgfmathsetmacro{\mf}{2*\mr+0.25*\mh}
			\pgfmathsetmacro{\md}{2*\mr+0.75*\mh}
			\pic[rotate=180] at (1.5*\mb,2*\mr+1.3*\mh){x_cos={0.4cm,0.2cm,ucbblue}};
			\node[ucbblue,anchor=south] at (1.5*\mb,2*\mr+1.3*\mh) {$x_1$};
			\pic at (\mb,0){wagen={\mh,\mb,\mr,$m$}};
			\pic {feder={0,\mb,\mf,$c$}};
			\pic {daempfer={0,\mb,\md,$d$}};
			% Wand
			\draw[](0,0) --++ (0,\mh+2*\mr);
			\fill[pattern={Lines[angle=45, distance={1pt}]},pattern color=black](0,0) rectangle (-3pt,\mh+2*\mr);
		}
	}
}
% ------------------ Zweimassenschwinger ------------------
% Pic: einmassenschwinger
% Parameter: #1 Höhe, #2 Breite,#3 Name der Masse, #4 Federsteifigkeit, #5 Dämpfung
\tikzset{
	pics/zweimassenschwinger/.style args={#1,#2,#3,#4,#5}{
		code={
			\pgfmathsetmacro{\mh}{0.8*#1}
			\pgfmathsetmacro{\mb}{0.5*#2}
			\pgfmathsetmacro{\mr}{0.1*#1}
			\pgfmathsetmacro{\mf}{2*\mr+0.25*\mh}
			\pgfmathsetmacro{\md}{2*\mr+0.75*\mh}
			% Wand
			\draw[](0,0) --++ (0,\mh+2*\mr);
			\fill[pattern={Lines[angle=45, distance={1pt}]},pattern color=black](0,0) rectangle (-3pt,\mh+2*\mr);
			% Schwinger
			\pic[rotate=180] at (1.5*\mb,2*\mr+1.3*\mh){x_cos={0.4cm,0.2cm,ucbblue}};
			\node[ucbblue,anchor=south] at (1.5*\mb,2*\mr+1.3*\mh) {$x_1$};
			\pic at (\mb,0){wagen={\mh,\mb,\mr,$m_1$}};
			\pic {feder={0,\mb,\mf,$c_1$}};
			\pic {daempfer={0,\mb,\md,$d_1$}};
			\begin{scope}[xshift=2*\mb cm]
				% Schwinger
				\pic[rotate=180] at (1.5*\mb,2*\mr+1.3*\mh){x_cos={0.4cm,0.2cm,ucbblue}};
				\node[ucbblue,anchor=south] at (1.5*\mb,2*\mr+1.3*\mh) {$x_2$};
				\pic at (\mb,0){wagen={\mh,\mb,\mr,$m_2$}};
				\pic {feder={0,\mb,\mf,$c_2$}};
				\pic {daempfer={0,\mb,\md,$d_2$}};
			\end{scope}
			\begin{scope}[xshift=4*\mb cm]
				\pic {feder={0,\mb,\mf,$c_3$}};
				\pic {daempfer={0,\mb,\md,$d_3$}};
			\end{scope}
			\begin{scope}[xshift=5*\mb cm]
				% Wand
				\draw[](0,0) --++ (0,\mh+2*\mr);
				\fill[pattern={Lines[angle=-45, distance={1pt}]},pattern color=black](0,0) rectangle (3pt,\mh+2*\mr);
			\end{scope}
		}
	}
}
% ------------------ Zweimassenschwinger m-k ------------------
% Pic: einmassenschwinger
% Parameter: #1 Höhe, #2 Breite,#3 Name der Masse, #4 Federsteifigkeit, #5 Dämpfung
\tikzset{
	pics/zweimassenschwinger_mk/.style args={#1,#2,#3,#4}{
		code={
			\pgfmathsetmacro{\mh}{0.8*#1}
			\pgfmathsetmacro{\mb}{0.5*#2}
			\pgfmathsetmacro{\mr}{0.1*#1}
			\pgfmathsetmacro{\mf}{2*\mr+0.5*\mh}
			% Wand
			\draw[](0,0) --++ (0,\mh+2*\mr);
			\fill[pattern={Lines[angle=45, distance={1pt}]},pattern color=black](0,0) rectangle (-3pt,\mh+2*\mr);
			% Schwinger
			\pic[rotate=180] at (1.5*\mb,2*\mr+1.3*\mh){x_cos={0.4cm,0.2cm,ucbblue}};
			\node[ucbblue,anchor=south] at (1.5*\mb,2*\mr+1.3*\mh) {$x_1$};
			\pic at (\mb,0){wagen={\mh,\mb,\mr,$m_1$}};
			\pic {feder={0,\mb,\mf,$c_1$}};
			\pic {daempfer={0,\mb,\md,$d_1$}};
			\begin{scope}[xshift=2*\mb cm]
				% Schwinger
				\pic[rotate=180] at (1.5*\mb,2*\mr+1.3*\mh){x_cos={0.4cm,0.2cm,ucbblue}};
				\node[ucbblue,anchor=south] at (1.5*\mb,2*\mr+1.3*\mh) {$x_2$};
				\pic at (\mb,0){wagen={\mh,\mb,\mr,$m_2$}};
				\pic {feder={0,\mb,\mf,$c_2$}};
			\end{scope}
			\begin{scope}[xshift=4*\mb cm]
				\pic {feder={0,\mb,\mf,$c_3$}};
			\end{scope}
			\begin{scope}[xshift=5*\mb cm]
				% Wand
				\draw[](0,0) --++ (0,\mh+2*\mr);
				\fill[pattern={Lines[angle=-45, distance={1pt}]},pattern color=black](0,0) rectangle (3pt,\mh+2*\mr);
			\end{scope}
		}
	}
}
% ------------------ Zweimassenschwinger m-k mode  ------------------
% Pic: einmassenschwinger
% Parameter: #1 Höhe, #2 Breite,#3 Name der Masse, #4 Federsteifigkeit, #5 Dämpfung
\tikzset{
	pics/zweimassenschwinger_mk_mode/.style args={#1,#2,#3,#4}{
		code={
			\pgfmathsetmacro{\mh}{0.8*#1}
			\pgfmathsetmacro{\mb}{0.5*#2}
			\pgfmathsetmacro{\mr}{0.1*#1}
			\pgfmathsetmacro{\mf}{2*\mr+0.5*\mh}
			% Schwinger
			\pic[rotate=180] at (1.5*\mb,2*\mr+1.3*\mh){x_cos={0.4cm,0.2cm,ucbblue}};
			\node[ucbblue,anchor=south] at (1.5*\mb,2*\mr+1.3*\mh) {$x_1$};
			\pic at (\mb,0){wagen_ob={\mh,\mb,\mr,$m_1$}};
			\pic {feder={0,\mb,\mf,$$}};
			\begin{scope}[xshift=2*\mb cm]
				% Schwinger
				\pic[rotate=180] at (1.5*\mb,2*\mr+1.3*\mh){x_cos={0.4cm,0.2cm,ucbblue}};
				\node[ucbblue,anchor=south] at (1.5*\mb,2*\mr+1.3*\mh) {$x_2$};
				\pic at (\mb,0){wagen_ob={\mh,\mb,\mr,$m_2$}};
				\pic {feder={0,\mb,\mf,$$}};
			\end{scope}
			\begin{scope}[xshift=4*\mb cm]
				\pic {feder={0,\mb,\mf,$$}};
			\end{scope}
			\begin{scope}[on background layer]
				\foreach \x in {0,0.25,...,5}{
					\draw[gray!30, thin] (\x,1.3) -- (\x,-6.);
				}
			\end{scope}
			\foreach \ys in {0,1.5,...,6}{
				\begin{scope}[yshift=-\ys cm]
					\draw[](0,\mh+2*\mr) -- (0,0) -- (5*\mb,0) --++ (0,\mh+2*\mr);
					\fill[pattern={Lines[angle=45, distance={1pt}]},pattern color=black](0,0) rectangle (-3pt,\mh+2*\mr);
					\fill[pattern={Lines[angle=45, distance={1pt}]},pattern color=black](-3pt,0) rectangle (5*\mb cm+3pt,-\mr);
					\fill[pattern={Lines[angle=45, distance={1pt}]},pattern color=black] (5*\mb cm,0) rectangle (5*\mb cm+3pt,\mh+2*\mr);
			\end{scope}
		}
		}
	}
}
% ------------------ Zweimassenschwinger m-k ------------------
% Pic: einmassenschwinger
% Parameter: #1 Höhe, #2 Breite,#3 Name der Masse, #4 Federsteifigkeit, #5 Dämpfung
\tikzset{
	pics/zweimassenschwinger_mdk_nicht_durchdringend/.style args={#1,#2,#3,#4,#5}{
		code={
			\pgfmathsetmacro{\mh}{0.8*#1}
			\pgfmathsetmacro{\mb}{0.5*#2}
			\pgfmathsetmacro{\mr}{0.1*#1}
			\pgfmathsetmacro{\mf}{2*\mr+0.25*\mh}
			\pgfmathsetmacro{\mfs}{2*\mr+0.5*\mh}
			\pgfmathsetmacro{\md}{2*\mr+0.75*\mh}
			% Wand
			\draw[](0,0) --++ (0,\mh+2*\mr);
			\fill[pattern={Lines[angle=45, distance={1pt}]},pattern color=black](0,0) rectangle (-3pt,\mh+2*\mr);
			% Schwinger
			\pic[rotate=180] at (1.5*\mb,2*\mr+1.3*\mh){x_cos={0.4cm,0.2cm,ucbblue}};
			\node[ucbblue,anchor=south] at (1.5*\mb,2*\mr+1.3*\mh) {$x_1$};
			\pic at (\mb,0){wagen={\mh,\mb,\mr,$m_1$}};
			\pic {feder={0,\mb,\mfs,$c_1$}};
			\begin{scope}[xshift=2*\mb cm]
				% Schwinger
				\pic[rotate=180] at (1.5*\mb,2*\mr+1.3*\mh){x_cos={0.4cm,0.2cm,ucbblue}};
				\node[ucbblue,anchor=south] at (1.5*\mb,2*\mr+1.3*\mh) {$x_2$};
				\pic at (\mb,0){wagen={\mh,\mb,\mr,$m_2$}};
				\pic {feder={0,\mb,\mf,$c_2$}};
				\pic {daempfer={0,\mb,\md,$d_2$}};
			\end{scope}
			\begin{scope}[xshift=4*\mb cm]
				\pic {feder={0,\mb,\mfs,$c_3$}};
			\end{scope}
			\begin{scope}[xshift=5*\mb cm]
				% Wand
				\draw[](0,0) --++ (0,\mh+2*\mr);
				\fill[pattern={Lines[angle=-45, distance={1pt}]},pattern color=black](0,0) rectangle (3pt,\mh+2*\mr);
			\end{scope}
		}
	}
}
% ------------------ Zweimassenschwinger m-k ------------------
% Pic: einmassenschwinger
% Parameter: #1 Höhe, #2 Breite,#3 Name der Masse, #4 Federsteifigkeit, #5 Dämpfung
\tikzset{
	pics/zweimassenschwinger_mdk_rayleigh_1/.style args={#1,#2,#3,#4,#5}{
		code={
			\pgfmathsetmacro{\mh}{0.8*#1}
			\pgfmathsetmacro{\mb}{0.5*#2}
			\pgfmathsetmacro{\mr}{0.1*#1}
			\pgfmathsetmacro{\mf}{2*\mr+0.25*\mh}
			\pgfmathsetmacro{\mfs}{2*\mr+0.5*\mh}
			\pgfmathsetmacro{\md}{2*\mr+0.75*\mh}
			% Wand
			\draw[](0,0) --++ (0,\mh+2*\mr);
			\fill[pattern={Lines[angle=45, distance={1pt}]},pattern color=black](0,0) rectangle (-3pt,\mh+2*\mr);
			% Schwinger
			\pic[rotate=180] at (1.5*\mb,2*\mr+1.3*\mh){x_cos={0.4cm,0.2cm,ucbblue}};
			\node[ucbblue,anchor=south] at (1.5*\mb,2*\mr+1.3*\mh) {$x_1$};
			\pic at (\mb,0){wagen={\mh,\mb,\mr,$m$}};
			\pic {feder={0,\mb,\mf,$c$}};
			\pic {daempfer={0,\mb,\md,$d$}};
			\begin{scope}[xshift=2*\mb cm]
				% Schwinger
				\pic[rotate=180] at (1.5*\mb,2*\mr+1.3*\mh){x_cos={0.4cm,0.2cm,ucbblue}};
				\node[ucbblue,anchor=south] at (1.5*\mb,2*\mr+1.3*\mh) {$x_2$};
				\pic at (\mb,0){wagen={\mh,\mb,\mr,$m$}};
				\pic {feder={0,\mb,\mfs,$c$}};
			\end{scope}
			\begin{scope}[xshift=4*\mb cm]
				\pic {feder={0,\mb,\mf,$c$}};
				\pic {daempfer={0,\mb,\md,$d$}};
			\end{scope}
			\begin{scope}[xshift=5*\mb cm]
				% Wand
				\draw[](0,0) --++ (0,\mh+2*\mr);
				\fill[pattern={Lines[angle=-45, distance={1pt}]},pattern color=black](0,0) rectangle (3pt,\mh+2*\mr);
			\end{scope}
		}
	}
}
% ------------------ Zweimassenschwinger m-k ------------------
% Pic: einmassenschwinger
% Parameter: #1 Höhe, #2 Breite,#3 Name der Masse, #4 Federsteifigkeit, #5 Dämpfung
\tikzset{
	pics/zweimassenschwinger_mdk_rayleigh_2/.style args={#1,#2,#3,#4,#5}{
		code={
			\pgfmathsetmacro{\mh}{0.8*#1}
			\pgfmathsetmacro{\mb}{0.5*#2}
			\pgfmathsetmacro{\mr}{0.1*#1}
			\pgfmathsetmacro{\mf}{2*\mr+0.25*\mh}
			\pgfmathsetmacro{\mfs}{2*\mr+0.5*\mh}
			\pgfmathsetmacro{\md}{2*\mr+0.75*\mh}
			% Wand
			\draw[](0,0) --++ (0,\mh+2*\mr);
			\fill[pattern={Lines[angle=45, distance={1pt}]},pattern color=black](0,0) rectangle (-3pt,\mh+2*\mr);
			% Schwinger
			\pic[rotate=180] at (1.5*\mb,2*\mr+1.3*\mh){x_cos={0.4cm,0.2cm,ucbblue}};
			\node[ucbblue,anchor=south] at (1.5*\mb,2*\mr+1.3*\mh) {$x_1$};
			\pic at (\mb,0){wagen={\mh,\mb,\mr,$m$}};
			\pic {feder={0,\mb,\mf,$c$}};
			\pic {daempfer={0,\mb,\md,$d$}};
			\begin{scope}[xshift=2*\mb cm]
				% Schwinger
				\pic[rotate=180] at (1.5*\mb,2*\mr+1.3*\mh){x_cos={0.4cm,0.2cm,ucbblue}};
				\node[ucbblue,anchor=south] at (1.5*\mb,2*\mr+1.3*\mh) {$x_2$};
				\pic at (\mb,0){wagen={\mh,\mb,\mr,$m$}};
				\pic {feder={0,\mb,\mfs,$c$}};
			\end{scope}
			\begin{scope}[xshift=4*\mb cm]
				\pic {feder={0,\mb,\mfs,$c$}};
			\end{scope}
			\begin{scope}[xshift=5*\mb cm]
				% Wand
				\draw[](0,0) --++ (0,\mh+2*\mr);
				\fill[pattern={Lines[angle=-45, distance={1pt}]},pattern color=black](0,0) rectangle (3pt,\mh+2*\mr);
			\end{scope}
		}
	}
}
% ------------------ Zweimassenschwinger m-k ------------------
% Pic: einmassenschwinger
% Parameter: #1 Höhe, #2 Breite,#3 Name der Masse, #4 Federsteifigkeit, #5 Dämpfung
\tikzset{
	pics/zweimassenschwinger_mdk_f_bsp1/.style args={#1,#2,#3,#4,#5}{
		code={
			\pgfmathsetmacro{\mh}{0.8*#1}
			\pgfmathsetmacro{\mb}{0.5*#2}
			\pgfmathsetmacro{\mr}{0.1*#1}
			\pgfmathsetmacro{\mf}{2*\mr+0.25*\mh}
			\pgfmathsetmacro{\mfs}{2*\mr+0.5*\mh}
			\pgfmathsetmacro{\md}{2*\mr+0.75*\mh}
			% Wand
			\draw[](0,0) --++ (0,\mh+2*\mr);
			\fill[pattern={Lines[angle=45, distance={1pt}]},pattern color=black](0,0) rectangle (-3pt,\mh+2*\mr);
			% Schwinger
			\pic[rotate=180] at (1.2*\mb,2*\mr+1.3*\mh){x_cos={0.4cm,0.2cm,ucbblue}};
			\node[ucbblue,anchor=south] at (1.2*\mb,2*\mr+1.3*\mh) {$x_1$};
			\pic at (\mb,0){wagen={\mh,\mb,\mr,$m$}};
			\pic {feder={0,\mb,\mf,$c$}};
			\pic {daempfer={0,\mb,\md,$d$}};
			\draw (1.5*\mb,2*\mr+1.*\mh)--+(0,2*\mr);
			\draw[red,-latex] (1.5*\mb,3*\mr+1.*\mh)--+(0.8,0)node[pos=.5,anchor=south]{$F_1$};
			\begin{scope}[xshift=2*\mb cm]
				% Schwinger
				\pic[rotate=180] at (1.2*\mb,2*\mr+1.3*\mh){x_cos={0.4cm,0.2cm,ucbblue}};
				\node[ucbblue,anchor=south] at (1.2*\mb,2*\mr+1.3*\mh) {$x_2$};
				\pic at (\mb,0){wagen={\mh,\mb,\mr,$m$}};
				\pic {feder={0,\mb,\mfs,$c$}};
				\draw (1.5*\mb,2*\mr+1.*\mh)--+(0,2*\mr);
				\draw[red,-latex] (1.5*\mb,3*\mr+1.*\mh)--+(0.8,0)node[pos=.5,anchor=south]{$F_2$};
			\end{scope}
			\begin{scope}[xshift=4*\mb cm]
				\pic {feder={0,\mb,\mfs,$c$}};
			\end{scope}
			\begin{scope}[xshift=5*\mb cm]
				% Wand
				\draw[](0,0) --++ (0,\mh+2*\mr);
				\fill[pattern={Lines[angle=-45, distance={1pt}]},pattern color=black](0,0) rectangle (3pt,\mh+2*\mr);
			\end{scope}
		}
	}
}
% ------------------ Zweimassenschwinger m-k ------------------
% Pic: einmassenschwinger
% Parameter: #1 Höhe, #2 Breite,#3 Name der Masse, #4 Federsteifigkeit, #5 Dämpfung
\tikzset{
	pics/zweimassenschwinger_mk_f/.style args={#1,#2,#3,#4}{
		code={
			\pgfmathsetmacro{\mh}{0.8*#1}
			\pgfmathsetmacro{\mb}{0.5*#2}
			\pgfmathsetmacro{\mr}{0.1*#1}
			\pgfmathsetmacro{\mf}{2*\mr+0.25*\mh}
			\pgfmathsetmacro{\mfs}{2*\mr+0.5*\mh}
			\pgfmathsetmacro{\md}{2*\mr+0.75*\mh}
			% Wand
			\draw[](0,0) --++ (0,\mh+2*\mr);
			\fill[pattern={Lines[angle=45, distance={1pt}]},pattern color=black](0,0) rectangle (-3pt,\mh+2*\mr);
			% Schwinger
			\pic[rotate=180] at (1.2*\mb,2*\mr+1.3*\mh){x_cos={0.4cm,0.2cm,ucbblue}};
			\node[ucbblue,anchor=south] at (1.2*\mb,2*\mr+1.3*\mh) {$x_1$};
			\pic at (\mb,0){wagen={\mh,\mb,\mr,$m$}};
			\pic {feder={0,\mb,\mfs,$c$}};
			\draw (1.5*\mb,2*\mr+1.*\mh)--+(0,2*\mr);
			\draw[red,-latex] (1.5*\mb,3*\mr+1.*\mh)--+(0.8,0)node[pos=.5,anchor=south]{$F_1$};
			\begin{scope}[xshift=2*\mb cm]
				% Schwinger
				\pic[rotate=180] at (1.2*\mb,2*\mr+1.3*\mh){x_cos={0.4cm,0.2cm,ucbblue}};
				\node[ucbblue,anchor=south] at (1.2*\mb,2*\mr+1.3*\mh) {$x_2$};
				\pic at (\mb,0){wagen={\mh,\mb,\mr,$m$}};
				\pic {feder={0,\mb,\mfs,$c$}};
				\draw (1.5*\mb,2*\mr+1.*\mh)--+(0,2*\mr);
				\draw[red,-latex] (1.5*\mb,3*\mr+1.*\mh)--+(0.8,0)node[pos=.5,anchor=south]{$F_2$};
			\end{scope}
			\begin{scope}[xshift=4*\mb cm]
				\pic {feder={0,\mb,\mfs,$c$}};
			\end{scope}
			\begin{scope}[xshift=5*\mb cm]
				% Wand
				\draw[](0,0) --++ (0,\mh+2*\mr);
				\fill[pattern={Lines[angle=-45, distance={1pt}]},pattern color=black](0,0) rectangle (3pt,\mh+2*\mr);
			\end{scope}
		}
	}
}
% ------------------ Dreimassenschwinger m-k ------------------
% Pic: einmassenschwinger
% Parameter: #1 Höhe, #2 Breite,#3 Name der Masse, #4 Federsteifigkeit, #5 Dämpfung
\tikzset{
	pics/dreimassenschwinger_mdk_durchdringend/.style args={#1,#2,#3,#4,#5,#6}{
		code={
			\pgfmathsetmacro{\mh}{0.8*#1}
			\pgfmathsetmacro{\mb}{0.5*#2}
			\pgfmathsetmacro{\mr}{0.1*#1}
			\pgfmathsetmacro{\mf}{2*\mr+0.25*\mh}
			\pgfmathsetmacro{\mfs}{2*\mr+0.5*\mh}
			\pgfmathsetmacro{\md}{2*\mr+0.75*\mh}
			% Wand
			\draw[](0,0) --++ (0,\mh+2*\mr);
			\fill[pattern={Lines[angle=45, distance={1pt}]},pattern color=black](0,0) rectangle (-3pt,\mh+2*\mr);
			% Schwinger
			\pic[rotate=180] at (1.5*\mb,2*\mr+1.3*\mh){x_cos={0.4cm,0.2cm,ucbblue}};
			\node[ucbblue,anchor=south] at (1.5*\mb,2*\mr+1.3*\mh) {$x_1$};
			\pic at (\mb,0){wagen={\mh,\mb,\mr,$m_1$}};
			\pic {feder={0,\mb,\mf,$c_1$}};
			\pic {daempfer={0,\mb,\md,$d_1$}};
			\begin{scope}[xshift=2*\mb cm]
				% Schwinger
				\pic[rotate=180] at (1.5*\mb,2*\mr+1.3*\mh){x_cos={0.4cm,0.2cm,ucbblue}};
				\node[ucbblue,anchor=south] at (1.5*\mb,2*\mr+1.3*\mh) {$x_2$};
				\pic at (\mb,0){wagen={\mh,\mb,\mr,$m_2$}};
				\pic {feder={0,\mb,\mfs,$c_2$}};
			\end{scope}
			\begin{scope}[xshift=4*\mb cm]
				% Schwinger
				\pic[rotate=180] at (1.5*\mb,2*\mr+1.3*\mh){x_cos={0.4cm,0.2cm,ucbblue}};
				\node[ucbblue,anchor=south] at (1.5*\mb,2*\mr+1.3*\mh) {$x_3$};
				\pic at (\mb,0){wagen={\mh,\mb,\mr,$m_3$}};
				\pic {feder={0,\mb,\mfs,$c_3$}};
			\end{scope}
			\begin{scope}[xshift=6*\mb cm]
				\pic {feder={0,\mb,\mfs,$c_4$}};
			\end{scope}
			\begin{scope}[xshift=7*\mb cm]
				% Wand
				\draw[](0,0) --++ (0,\mh+2*\mr);
				\fill[pattern={Lines[angle=-45, distance={1pt}]},pattern color=black](0,0) rectangle (3pt,\mh+2*\mr);
			\end{scope}
		}
	}
}
% ------------------ n-massenschwinger ------------------
% Pic: einmassenschwinger
% Parameter: #1 Höhe, #2 Breite,#3 Name der Masse, #4 Federsteifigkeit, #5 Dämpfung
\tikzset{
	pics/nmassenschwinger/.style args={#1,#2,#3,#4,#5}{
		code={
			\pgfmathsetmacro{\mh}{0.8*#1}
			\pgfmathsetmacro{\mb}{0.5*#2}
			\pgfmathsetmacro{\mr}{0.1*#1}
			\pgfmathsetmacro{\mf}{2*\mr+0.25*\mh}
			\pgfmathsetmacro{\md}{2*\mr+0.75*\mh}
			% Wand
			\draw[](0,0) --++ (0,\mh+2*\mr);
			\fill[pattern={Lines[angle=45, distance={1pt}]},pattern color=black](0,0) rectangle (-3pt,\mh+2*\mr);
			% Schwinger
			\pic at (\mb,0){wagen={\mh,\mb,\mr,$m_1$}};
			\pic {feder={0,\mb,\mf,$c_1$}};
			\pic {daempfer={0,\mb,\md,$d_1$}};
			\begin{scope}[xshift=2*\mb cm]
				% Schwinger
				\pic {feder={0,\mb,\mf,$c_2$}};
				\pic {daempfer={0,\mb,\md,$d_2$}};
			\end{scope}
			\begin{scope}[xshift=5*\mb cm]
				\pic at (\mb,0){wagen={\mh,\mb,\mr,$m_n$}};
				\pic {feder={0,\mb,\mf,$c_{n}$}};
				\pic {daempfer={0,\mb,\md,$d_{n}$}};
			\end{scope}
			\begin{scope}[xshift=7*\mb cm]
				\pic {feder={0,\mb,\mf,$c_{n+1}$}};
				\pic {daempfer={0,\mb,\md,$d_{n+1}$}};
			\end{scope}
			\begin{scope}[xshift=8*\mb cm]
				% Wand
				\draw[](0,0) --++ (0,\mh+2*\mr);
				\fill[pattern={Lines[angle=-45, distance={1pt}]},pattern color=black](0,0) rectangle (3pt,\mh+2*\mr);
			\end{scope}
		}
	}
}
